\documentclass[
	12pt,				% tamanho da fonte
	openright,			% capítulos começam em pág ímpar (insere página vazia caso preciso)
	oneside,			% para impressão em apenas um lado. Substitua por twoside para frente e verso
	a4paper,			% tamanho do papel. 
	english,			% idioma adicional para hifenização
	brazil				% o último idioma é o principal do documento
	]{abntex2}

% --- PACOTES BÁSICOS ---
\usepackage{lmodern}			% Usa a fonte Latin Modern			
\usepackage[T1]{fontenc}		% Selecao de codigos de fonte.
\usepackage[utf8]{inputenc}		% Codificacao do documento (conversão automática dos acentos)
\usepackage{lastpage}			% Usado no Cabeçalho
\usepackage{indentfirst}		% Indenta o primeiro parágrafo de cada seção.
\usepackage{color}				% Controle das cores
\usepackage{graphicx}			% Inclusão de gráficos
\usepackage{microtype} 			% para melhorias de justificação
\usepackage{tabularx}
\usepackage{booktabs}

% --- PACOTES DE CITAÇÕES ---
\usepackage[brazilian,hyperpageref]{backref}	 % Paginas com as citações na bibl
\usepackage[alf]{abntex2cite}	% Citações padrão ABNT

% --- CONFIGURAÇÕES DE APARÊNCIA ---
\definecolor{blue}{RGB}{41,5,195}
\hypersetup{
    pdftitle={Dossiê Estratégico IntraClinica}, 
    pdfauthor={Bruno},
    pdfsubject={Soberania Digital e Inteligência Clínica},
    pdfkeywords={Gestão Clínica, Soberania de Dados, Experiência Clínica Soberana (Ecossistema IntraClinica), IA Local},
    colorlinks=true,
    linkcolor=blue,
    citecolor=blue,
    filecolor=magenta,
    urlcolor=blue,
    bookmarksdepth=4
}

% --- DADOS DO TRABALHO ---
\titulo{Dossiê Estratégico IntraClinica:\\Soberania Digital e Excelência na Prática Médica (Experiência Clínica Soberana (Ecossistema IntraClinica))}
\autor{Ecossistema IntraClinica}
\local{São José do Rio Preto - SP}
\data{2026}
\instituicao{IntraClinica - Inovação e Curadoria Médica}
\tipotrabalho{Relatório Técnico}
\preambulo{Relatório técnico-estratégico apresentado à Diretoria Médica para fins de expansão de rede e definição de protocolos de curadoria para médicos parceiros.}

% --- CONFIGURAÇÕES DE PÁGINA ---
\makeatletter
\setlength{\parindent}{1.3cm}
\setlength{\parskip}{0.2cm}  
\makeatother

% --- INÍCIO DO DOCUMENTO ---
\begin{document}

\selectlanguage{brazil}
\frenchspacing 

% --- CAPA ---
\imprimircapa

% --- FOLHA DE ROSTO ---
\imprimirfolhaderosto

% --- RESUMO ---
\begin{resumo}
 O IntraClinica é um ecossistema de gestão clínica desenhado para priorizar a soberania de dados e a experiência do profissional de saúde (Experiência Clínica Soberana (Ecossistema IntraClinica)). Este dossiê detalha a arquitetura do sistema, fundamentada em tecnologias de código aberto e persistência local, visando mitigar a dependência de plataformas proprietárias. São abordados fluxos de uso-case, desde a recepção até o atendimento clínico monomanual, e a gestão inteligente de insumos via inteligência artificial local.
 
 \vspace{\onelineskip}
 \noindent
 \textbf{Palavras-chave}: Soberania de Dados. Gestão Clínica. Experiência Clínica Soberana (Ecossistema IntraClinica). Inteligência Artificial.
\end{resumo}

% --- SUMÁRIO ---
\pdfbookmark[0]{\contentsname}{toc}
\tableofcontents*
\cleardoublepage

% --- ELEMENTOS TEXTUAIS ---
\textual

\chapter{Introdução}

O cenário atual da gestão de saúde no Brasil é marcado por uma dicotomia: de um lado, sistemas legados robustos, porém burocráticos e de difícil usabilidade; de outro, soluções modernas em nuvem que, embora ágeis, sequestram a soberania dos dados do profissional médico e impõem custos operacionais crescentes.

O projeto \textbf{IntraClinica} surge como uma terceira via. Fundamentado no binômio \textit{``Medicine for Humans. Intelligence for Business.''}, o sistema busca devolver ao médico a autoridade sobre o seu fluxo de trabalho e sobre a informação gerada em sua prática.

\section{Objetivos}

Este relatório tem como objetivo principal detalhar a arquitetura e os processos do IntraClinica para a Diretoria Médica e parceiros estratégicos. Os pilares deste desenvolvimento são:

\begin{itemize}
    \item \textbf{Soberania Radical:} Garantir que a posse física e lógica dos dados resida onde o médico exerce sua vontade.
    \item \textbf{Excelência em Experiência Médica (Experiência Clínica Soberana (Ecossistema IntraClinica)):} Desenvolver interfaces que não concorram pela atenção do médico com o paciente.
    \item \textbf{Resiliência Operacional:} Implementar infraestrutura capaz de operar independente de instabilidades em serviços externos.
\end{itemize}

\section{Justificativa}

A necessidade de um sistema como o IntraClinica foi identificada através do diagnóstico de ``hemorragia financeira'' em clínicas de médio porte, onde o descontrole sobre o almoxarifado e a fadiga gerada por softwares hostis resultam em prejuízos diretos e queda na qualidade do atendimento clínico.

\chapter{Experiência Clínica Soberana}

O ecossistema IntraClinica é fundamentado na busca pela fluidez operacional e na soberania técnica do profissional de saúde. Diferente das arquiteturas tradicionais de mercado, o design do sistema considera as condições específicas do ambiente clínico de alto desempenho: assepsia, ergonomia e o valor primordial do contato visual com o paciente.

\section{Otimização de Processos e Prevenção de Perdas}

A concepção da plataforma baseou-se na identificação de gargalos operacionais comuns em unidades de saúde de médio porte, visando elevar o padrão de eficiência para além das soluções de referência atuais:

\subsection{Eficiência na Gestão de Insumos}
Em ecossistemas complexos, a gestão de estoques críticos pode sofrer com a falta de rastreabilidade instantânea. O IntraClinica propõe um modelo de baixa atômica que automatiza o registro de consumo, garantindo a previsibilidade financeira e a segurança do paciente através de dados em tempo real.

\subsection{Simplificação de Fluxo e Redução de Fadiga Digital}
Sistemas legados tendem a apresentar interfaces saturadas de informações secundárias que podem gerar fadiga cognitiva no profissional de saúde. O IntraClinica prioriza a agilidade do atendimento, entregando uma interface limpa que atua como um suporte silencioso à decisão clínica.

\section{Interface Monomanual e a Preservação da Atenção}

O IntraClinica implementa a filosofia de que o software deve ser uma extensão natural da vontade médica. Para tal, as interfaces foram otimizadas para uso em tablets com operação monomanual, permitindo que o médico registre anamneses e evoluções sem comprometer a ergonomia ou a interação humana com o paciente.

\chapter{Fundação Tecnológica e Soberania}

A arquitetura do IntraClinica é dividida em camadas que garantem a integridade do dado e a soberania do profissional de saúde.

\section{Persistência e Backend-as-a-Service (BaaS)}

O sistema utiliza o \textbf{Supabase (PostgreSQL 16)} como fonte única da verdade. A escolha pelo PostgreSQL garante que o dado seja relacional, auditável e exportável a qualquer momento, sem lock-in de fornecedor.

\subsection{Isolamento Multi-Tenancy via RLS}
Para a expansão da rede de parceiros, a segurança é garantida no nível do banco de dados através de \textbf{Row Level Security (RLS)}. Cada registro clínico ou de estoque é vinculado a um identificador único de clínica (\textit{clinic\_id}). O motor do Postgres garante que um usuário só acesse dados para os quais seu perfil (\textit{profiles}) tenha autorização explícita, impedindo vazamentos transversais de dados.

\section{Modelo de Dados Relacional}

O coração do IntraClinica é um modelo relacional normalizado, conforme ilustrado graficamente no diagrama \texttt{db\_schema.svg} (ver Apêndice A). Esta estrutura foi desenhada para garantir a integridade dos dados clínicos e a rastreabilidade total do inventário. A estrutura é centrada no conceito de \textbf{Multi-Tenancy}, onde cada tabela core possui uma chave estrangeira para a entidade \texttt{clinic}.

\subsection{Entidades e Relacionamentos (ER)}

A modelagem utiliza o padrão de \textbf{Supertipo/Subtipo} para a entidade \texttt{Actor}. Um \texttt{Actor} representa qualquer ente humano no sistema, que pode se especializar como um \texttt{Patient} (paciente) ou um \texttt{User} (profissional/usuário).

Principais módulos do banco:
\begin{itemize}
    \item \textbf{Módulo de Identidade:} Tabelas \texttt{clinic}, \texttt{actor}, \texttt{user} e \texttt{patient}.
    \item \textbf{Módulo Clínico:} Tabelas \texttt{appointment} (agendamentos) e \texttt{clinical\_record} (prontuário).
    \item \textbf{Módulo de Suprimentos:} Tabelas \texttt{inventory\_item}, \texttt{procedure\_recipe} e \texttt{inventory\_movement}.
    \item \textbf{Módulo Financeiro:} Tabelas \texttt{financial\_transaction} e \texttt{financial\_category}.
\end{itemize}

O diagrama detalhado deste modelo, representando a genealogia das tabelas e suas restrições, encontra-se no \textbf{Apêndice A}.

\section{Reatividade Granular via Angular 21}

O frontend é desenvolvido em \textbf{Angular 21}, utilizando a nova arquitetura de \textbf{Signals}. Isso permite que a interface reaja instantaneamente a mudanças no banco de dados sem a necessidade de recarregamentos pesados de página, otimizando o fluxo de trabalho na recepção e no consultório.

\section{Infraestrutura de Engenharia e Governança (Axio Nexus)}

O desenvolvimento e a manutenção do IntraClinica não dependem de supervisão humana constante para integridade de dados, mas são governados pela malha neural \textbf{Axio Nexus}. Diferente de um módulo do sistema, o Nexus atua como a infraestrutura de inteligência que sustenta o projeto.

O motor \textbf{Weaver Worker} realiza a mineração contínua do código-fonte e da documentação técnica, operando em três níveis:
\begin{itemize}
    \item \textbf{Auditoria de Conformidade:} Verifica se as implementações em Angular e PostgreSQL respeitam as premissas de soberania e Experiência Clínica Soberana (Ecossistema IntraClinica).
    \item \textbf{Tecelagem Topológica (T4):} Conecta decisões de arquitetura à implementação real, prevenindo a dívida técnica e a perda de genealogia do projeto.
    \item \textbf{Persistência Cognitiva:} Garante que a inteligência acumulada durante o desenvolvimento seja soberana e local, independente de provedores de IA externos.
\end{itemize}

\chapter{Fluxos Operacionais e Casos de Uso}

Este capítulo descreve os fluxos vitais do sistema e como eles materializam a visão Experiência Clínica Soberana (Ecossistema IntraClinica) no dia a dia da clínica.

\section{Recepção e Triagem de Alta Performance}

O fluxo de entrada foi desenhado para ser o mais silencioso possível.
\begin{itemize}
    \item \textbf{Fluxo:} O secretariado realiza o check-in rápido. Os dados demográficos são validados e o paciente é inserido na fila de espera inteligente.
    \item \textbf{Impacto:} Quando o médico abre o prontuário, o sistema já apresenta o contexto demográfico e o histórico de retornos, permitindo que a consulta inicie sem perguntas repetitivas e burocráticas.
\end{itemize}

\section{O Ato Médico e a Máquina de Estados}

O status da consulta é gerido por uma máquina de estados rigorosa, impedindo saltos de etapa que poderiam comprometer a auditoria.
\begin{enumerate}
    \item \textbf{Agendado:} O registro inicial.
    \item \textbf{Aguardando:} O paciente está presente fisicamente.
    \item \textbf{Em Atendimento:} O médico iniciou a anamnese.
    \item \textbf{Realizado:} O atendimento foi concluído com a geração de documentos (receitas, atestados) ou registro de procedimentos.
\end{enumerate}

Essa progressão garante que o sistema capture o tempo real de cada fase do atendimento, gerando indicadores de performance (KPIs) precisos para a gestão da unidade.

\section{Comunicação e Engajamento Ético}

O sistema inclui um módulo de inteligência criativa que auxilia a clínica na produção de conteúdo educativo para redes sociais e lembretes de retorno via canais de mensagem, fortalecendo a autoridade médica sem violar os limites éticos da profissão.

\chapter{Governança: Estoque e Matriz de Acesso}

A sustentabilidade financeira e a segurança jurídica da rede IntraClinica residem no controle rigoroso sobre os insumos e sobre quem manipula a informação.

\section{O Guardião de Estoque Inteligente}

O estoque no IntraClinica não é um módulo isolado, mas uma consequência do atendimento.
\begin{itemize}
    \item \textbf{Receita de Procedimento:} Para cada ato médico (ex: aplicação de injetável, pequena cirurgia), existe uma ``receita'' cadastrada no sistema.
    \item \textbf{Baixa Atômica:} Ao finalizar o procedimento no prontuário, o sistema automaticamente propõe a baixa dos insumos associados. Se uma sutura foi realizada, o kit de sutura e os fios são deduzidos do inventário instantaneamente.
    \item \textbf{Alertas de Anomalia:} A inteligência local monitora desvios estatísticos. Se o consumo de um determinado profissional ou turno diverge do padrão para aquele procedimento, a Diretoria Médica recebe um alerta silencioso de auditoria.
\end{itemize}

\section{Gestão de Identidade e Acesso (IAM)}

Para viabilizar a curadoria de médicos parceiros, implementamos uma matriz de acesso baseada no princípio do menor privilégio.

\begin{table}[h!]
\centering
\caption{Matriz de Funções e Escopos}
\vspace{0.3cm}
\begin{tabular}{@{}ll@{}}
\toprule
\textbf{Papel (Role)} & \textbf{Responsabilidade e Acesso} \\ \midrule
Diretoria Médica & Gestão total, auditoria e definição de protocolos. \\
Médico Especialista & Acesso pleno ao prontuário e atendimento clínico. \\
Secretariado & Foco em agendamento, check-in e dados cadastrais. \\
Gestor de Estoque & Controle de suprimentos, validades e custos. \\ \bottomrule
\end{tabular}
\end{table}

Essa estrutura garante que a intimidade do paciente (contida no prontuário médico) seja acessada apenas por profissionais autorizados, cumprindo rigorosamente as normativas do CFM e da LGPD.

\chapter{Conclusão e Próximos Passos}

O IntraClinica encerra sua fase de desenvolvimento core com uma fundação sólida, resiliente e soberana. O sistema prova que é possível aliar alta tecnologia com a simplicidade exigida pela prática médica real.

\section{Metodologia de Evolução}

A evolução do IntraClinica é pautada pelo **Rigor Axiomático**. Cada nova funcionalidade é submetida à validação da malha neural Nexus, que atua como o engenheiro chefe do ecossistema, garantindo que o crescimento da rede de parceiros não comprometa a resiliência e a soberania do dado clínico.

\vfill
\centering
\textit{Este relatório foi compilado e verificado pela infraestrutura de inteligência Axio Nexus v3.3.} 🧪🧬🏥


% --- ELEMENTOS PÓS-TEXTUAIS ---
\postextual

\apendices
\chapter{Diagramas de Lógica e Fluxo}

Este apêndice contém as representações visuais dos processos automatizados e da estrutura de dados do sistema IntraClinica.

\section{Diagrama de Entidade e Relacionamento (ER)}

A estrutura relacional do banco de dados PostgreSQL do IntraClinica está representada graficamente no diagrama abaixo.

\begin{figure}[h!]
    \centering
    \includegraphics[width=\textwidth]{db_schema_v2.mermaid}
    \caption{Diagrama de Entidade e Relacionamento - IntraClinica}
    \label{fig:db_schema}
\end{figure}

O diagrama detalha as relações entre as entidades de identidade (\texttt{clinic}, \texttt{actor}), o módulo clínico (\texttt{appointment}, \texttt{clinical\_record}) e o motor de suprimentos (\texttt{inventory\_item}, \texttt{procedure\_recipe}).

\section{Fluxo de Baixa de Estoque Atômica}

O diagrama a seguir detalha como o sistema lida com a transação de baixa de insumos durante um procedimento clínico:

\begin{enumerate}
    \item O Médico finaliza o procedimento na interface Angular.
    \item O sistema dispara uma chamada RPC (\textit{Remote Procedure Call}) para o Supabase.
    \item O Postgres inicia uma transação atômica.
    \item Busca-se a receita cadastrada para aquele tipo de procedimento (\texttt{procedure\_recipe}).
    \item Para cada item, registra-se um rastro de movimentação (\texttt{inventory\_movement}) e deduz-se a quantidade do saldo atual (\texttt{inventory\_item}).
    \item A transação é confirmada e o médico recebe o log de auditoria na tela.
\end{enumerate}

\section{Diagrama de Estados da Consulta}

Representação da progressão lógica do atendimento:

\texttt{Agendado [Check-in] -> Aguardando [Chamada] -> Em Atendimento [Finalização] -> Realizado}

\section{Estrutura de Governança Nexus}

O ecossistema é governado pela arquitetura Axio Nexus, que atua como o engenheiro chefe e garantidor da integridade sistêmica:
\begin{itemize}
    \item \textbf{Tier 1:} Fragmentos brutos do código e documentação.
    \item \textbf{Tier 2:} Essência lógica destilada para governança.
    \item \textbf{Tier 4 (Topologia):} Links de causalidade que sustentam a inteligência do projeto.
\end{itemize}


\end{document}
